\documentclass[aspectratio=169]{beamer}
\usepackage{iftex}
\ifPDFTeX
  \usepackage[T1]{fontenc}
  \usepackage{lmodern}
\fi
\usepackage[czech]{babel}
\usepackage{amsmath}
\usepackage{listings}
\usepackage{tikz}
\usetikzlibrary{arrows.meta,positioning,shapes.multipart,babel}

\newcommand{\CVUTTemplatePath}{../../../utils/presentation-template}
\usepackage{\CVUTTemplatePath/beamerthemeCVUT}
\setbeamertemplate{footline}[frame number]
\beamertemplatenavigationsymbolsempty

\newcommand{\LANG}{CZ}
\newcommand{\FRONTPAGE}{dark}
\title{Úvod do objektově orientovaného programování}
\subtitle{Základy programování\\Matouš Cejnek\\U12110, FS, ČVUT v Praze}
\author{}
\institute{}
\date{}
\renewcommand{\headlineA}{}
\renewcommand{\headlineB}{}
\renewcommand{\headlineC}{}

\definecolor{CodeBlue}{RGB}{38,83,125}
\definecolor{SoftGray}{RGB}{242,244,247}
\definecolor{CodeGreen}{RGB}{38,125,94}
\setbeamerfont{normal text}{size=\large}
\setbeamerfont{frametitle}{size=\LARGE}
\setbeamercolor{block title}{fg=white,bg=CodeBlue}
\setbeamercolor{block body}{fg=black,bg=SoftGray}
\newcommand{\term}[1]{\textcolor{CodeBlue}{\textbf{#1}}}
\lstdefinestyle{python}{language=Python,basicstyle=\ttfamily\large,
  keywordstyle=\color{CodeBlue}\bfseries,stringstyle=\color{CodeGreen},
  commentstyle=\color{gray},showstringspaces=false,columns=fullflexible,keepspaces=true}

\begin{document}

% 0:00--0:01
\begin{frame}
  \clearpage\thispagestyle{empty}\titlepage
\end{frame}

% 0:01--0:03
\begin{frame}{Co si dnes odnesete?}
  \begin{itemize}
    \item Rozlišit třídu, objekt a instanci.
    \item Vytvořit třídu s atributy, konstruktorem a metodami.
    \item Vysvětlit roli \texttt{self} a sledovat změny stavu objektu.
    \item Číst jednoduchý diagram tříd a sekvenční diagram.
    \item Použít dědičnost, přepsání metody a polymorfismus.
  \end{itemize}
\end{frame}

\begin{frame}{Na co navazujeme?}
  \begin{itemize}
    \item Funkce pojmenovávají chování a přijímají parametry.
    \item Seznamy a slovníky seskupují hodnoty.
    \item Objekt spojí \term{data a operace nad nimi} do jednoho celku.
  \end{itemize}
  \begin{block}{OOP je nástroj organizace}
    Ne každý program potřebuje vlastní třídy. Používáme je, když model zpřehlední.
  \end{block}
\end{frame}

% 0:03--0:09
\begin{frame}[fragile]{Samostatná data se mohou rozcházet}
\begin{lstlisting}[style=python]
names = ["Eva", "Petr"]
balances = [700, 1200]

balances[0] += 200
\end{lstlisting}
  \begin{itemize}
    \item Vazbu mezi jménem a zůstatkem drží jen stejný index.
    \item Operace nad účtem jsou rozptýlené po programu.
    \item Chceme mít data i povolené změny pohromadě.
  \end{itemize}
\end{frame}

\begin{frame}{Třída a objekt}
  \begin{description}
    \item[Třída] předpis společné struktury a chování, například \texttt{BankAccount}.
    \item[Objekt] konkrétní hodnota vytvořená podle třídy, například účet Evy.
    \item[Instance] jiné označení objektu dané třídy.
  \end{description}
  \begin{block}{Jedna třída, mnoho objektů}
    Každý objekt má vlastní stav, ale sdílí definici metod své třídy.
  \end{block}
\end{frame}

\begin{frame}[fragile]{První třída}
\begin{lstlisting}[style=python]
class BankAccount:
    def __init__(self, owner, balance=0):
        self.owner = owner
        self.balance = balance

    def deposit(self, amount):
        self.balance += amount
\end{lstlisting}
  Název třídy se podle konvence píše \texttt{CamelCase}; názvy metod jako funkce.
\end{frame}

\begin{frame}[fragile]{Vytvoření objektu a volání metody}
\begin{lstlisting}[style=python]
eva = BankAccount("Eva", 700)
petr = BankAccount("Petr")

eva.deposit(200)
print(eva.balance)    # 900
print(petr.balance)   # 0
\end{lstlisting}
  \begin{itemize}
    \item Volání třídy vytvoří nový objekt.
    \item Tečka zpřístupní atribut nebo metodu konkrétního objektu.
  \end{itemize}
\end{frame}

% 0:09--0:16
\begin{frame}{Atribut, metoda, stav}
  \begin{description}
    \item[Atribut] pojmenovaná hodnota patřící objektu: \texttt{eva.balance}.
    \item[Metoda] funkce definovaná ve třídě: \texttt{eva.deposit(200)}.
    \item[Stav] aktuální soubor hodnot atributů objektu.
  \end{description}
  Metoda může stav číst, změnit nebo z něj vypočítat výsledek.
\end{frame}

\begin{frame}[fragile]{Co dělá \texttt{\_\_init\_\_}?}
\begin{lstlisting}[style=python]
eva = BankAccount("Eva", 700)
\end{lstlisting}
  \begin{enumerate}
    \item Python vytvoří nový prázdný objekt.
    \item Automaticky zavolá \texttt{\_\_init\_\_(nový\_objekt, "Eva", 700)}.
    \item \texttt{\_\_init\_\_} nastaví počáteční stav objektu.
  \end{enumerate}
  \begin{alertblock}{Přesný pojem}
    \texttt{\_\_init\_\_} je inicializátor; v úvodním kurzu se často zjednodušeně nazývá konstruktor.
  \end{alertblock}
\end{frame}

\begin{frame}{Co znamená \texttt{self}?}
  \begin{block}{\texttt{self}}
    Odkaz na konkrétní objekt, se kterým právě probíhá volání metody.
  \end{block}
  \begin{itemize}
    \item \texttt{eva.deposit(200)} odpovídá \texttt{BankAccount.deposit(eva, 200)}.
    \item Python první argument doplní při volání přes objekt.
    \item Jméno \texttt{self} je konvence, kterou vždy dodržujeme.
  \end{itemize}
\end{frame}

\begin{frame}[fragile]{Parametr a atribut nejsou totéž}
\begin{lstlisting}[style=python]
def __init__(self, owner, balance=0):
    self.owner = owner
    self.balance = balance
\end{lstlisting}
  \begin{itemize}
    \item \texttt{owner} je lokální parametr tohoto volání.
    \item \texttt{self.owner} je atribut, který zůstane v objektu.
    \item Pravá strana se vyhodnotí a levá určuje, kam hodnotu uložit.
  \end{itemize}
\end{frame}

\begin{frame}[fragile]{Každý objekt má vlastní stav}
\begin{lstlisting}[style=python]
a = BankAccount("Eva", 700)
b = BankAccount("Petr", 700)
a.deposit(200)

print(a.balance, b.balance)  # 900 700
\end{lstlisting}
  Proměnné \texttt{a} a \texttt{b} odkazují na dva různé objekty.
\end{frame}

% 0:16--0:23
\begin{frame}[fragile]{Metoda může vracet hodnotu}
\begin{lstlisting}[style=python]
class BankAccount:
    # __init__ and deposit are unchanged

    def can_withdraw(self, amount):
        return amount <= self.balance

    def describe(self):
        return f"{self.owner}: {self.balance:.2f} CZK"
\end{lstlisting}
  Stejně jako u funkcí rozlišujeme výpis pomocí \texttt{print} a návrat hodnoty.
\end{frame}

\begin{frame}[fragile]{Metoda může hlídat platné změny}
\begin{lstlisting}[style=python,basicstyle=\ttfamily\small]
def withdraw(self, amount):
    if amount <= 0:
        return False
    if amount > self.balance:
        return False
    self.balance -= amount
    return True
\end{lstlisting}
  \begin{block}{Zapouzdření}
    Objekt nabízí operace, které chrání pravidla jeho stavu; uživatel nemusí znát vnitřní postup.
  \end{block}
\end{frame}

\begin{frame}[fragile]{Metoda může volat jinou metodu}
\begin{lstlisting}[style=python]
def transfer_to(self, other, amount):
    if not self.withdraw(amount):
        return False
    other.deposit(amount)
    return True

eva.transfer_to(petr, 300)
\end{lstlisting}
  \texttt{self} je odesílající účet; \texttt{other} odkazuje na jiný objekt.
\end{frame}

\begin{frame}{Objekt je také hodnota}
  \begin{itemize}
    \item Lze ho uložit do proměnné nebo seznamu.
    \item Lze ho předat jako argument funkci či metodě.
    \item Více proměnných může odkazovat na tentýž objekt.
  \end{itemize}
  \begin{alertblock}{Důsledek sdíleného odkazu}
    Změna objektu přes jeden odkaz je vidět i přes druhý odkaz na stejný objekt.
  \end{alertblock}
\end{frame}

\begin{frame}[fragile]{Kdy třída pomáhá?}
\begin{lstlisting}[style=python]
accounts = [
    BankAccount("Eva", 700),
    BankAccount("Petr", 1200),
    BankAccount("Jana", 300),
]

for account in accounts:
    print(account.describe())
\end{lstlisting}
  Třída dává všem prvkům stejné rozhraní a každému vlastní stav.
\end{frame}

% 0:23--0:30
\begin{frame}{UML diagram třídy}
  \centering
  \begin{tikzpicture}
    \node[draw,rectangle split,rectangle split parts=3,align=left,
      text width=62mm,fill=SoftGray] {
      \centering\textbf{BankAccount}
      \nodepart{second}
      + owner: str\\
      + balance: float
      \nodepart{third}
      + deposit(amount)\\
      + withdraw(amount): bool\\
      + describe(): str
    };
  \end{tikzpicture}
  \begin{itemize}
    \item Horní část: jméno třídy; prostřední: atributy; dolní: metody.
    \item \texttt{+} značí veřejnou část rozhraní, \texttt{-} neveřejný detail.
  \end{itemize}
\end{frame}

\begin{frame}{UML není kód}
  \begin{itemize}
    \item Diagram ukazuje strukturu a vztahy, ne přesný postup metod.
    \item Typ za dvojtečkou pomáhá čtenáři, i když jej Python nemusí vynucovat.
    \item Závorky rozlišují metodu od atributu.
    \item Diagram třídy popisuje předpis, nikoli aktuální hodnoty jednoho objektu.
  \end{itemize}
\end{frame}

\begin{frame}{Sekvenční diagram: pořadí volání}
  \centering
  \vspace{5mm}

  \begin{tikzpicture}[>=Latex,scale=.76,transform shape]
    \node (u) at (0,0) {program};
    \node (e) at (4,0) {\texttt{eva:BankAccount}};
    \node (p) at (8,0) {\texttt{petr:BankAccount}};
    \draw[dashed] (u)--(0,-4.1); \draw[dashed] (e)--(4,-4.1); \draw[dashed] (p)--(8,-4.1);
    \draw[->] (0,-.8)--node[above]{\texttt{transfer\_to(petr, 300)}}(4,-.8);
    \draw[->] (4,-1.7)--node[above]{\texttt{withdraw(300)}}(4,-2.4);
    \draw[->] (4,-3.0)--node[above]{\texttt{deposit(300)}}(8,-3.0);
    \draw[->,dashed] (8,-3.5)--node[above]{návrat}(4,-3.5);
  \end{tikzpicture}
  \par\vspace{2mm}

  Diagram zachycuje, kdo komu a v jakém pořadí posílá volání.
\end{frame}

\begin{frame}{Diagram tříd vs. sekvenční diagram}
  \begin{columns}[T]
    \begin{column}{.46\textwidth}
      \begin{block}{Diagram tříd}
        Statická struktura: třídy, atributy, metody a vztahy.
      \end{block}
    \end{column}
    \begin{column}{.46\textwidth}
      \begin{block}{Sekvenční diagram}
        Dynamický průběh: konkrétní objekty a pořadí volání v čase.
      \end{block}
    \end{column}
  \end{columns}
\end{frame}

% 0:30--0:38
\begin{frame}{Dědičnost: specializace existující třídy}
  \begin{block}{Dědičnost}
    Potomek přebírá chování rodiče a může ho rozšířit nebo upravit.
  \end{block}
  \begin{itemize}
    \item \texttt{Employee} je obecný zaměstnanec.
    \item \texttt{Doctor} je specializovaný zaměstnanec.
    \item Dědičnost vyjadřuje vztah „je druhem“ (\emph{is-a}).
  \end{itemize}
\end{frame}

\begin{frame}[fragile]{Základní zápis dědičnosti}
\begin{lstlisting}[style=python]
class Employee:
    def describe(self):
        return "employee"

class Doctor(Employee):
    def examine(self, patient):
        return f"Examining {patient}"

d = Doctor()
print(d.describe())       # inherited method
\end{lstlisting}
\end{frame}

\begin{frame}[fragile]{Přepsání metody (override)}
\begin{lstlisting}[style=python]
class Employee:
    def describe(self):
        return "employee"

class Doctor(Employee):
    def describe(self):
        return "doctor"

print(Doctor().describe())  # doctor
\end{lstlisting}
  Potomek definuje metodu stejného jména; při volání se použije jeho verze.
\end{frame}

\begin{frame}[fragile]{Vlastní \texttt{\_\_init\_\_} přepíše rodičovský}
\begin{lstlisting}[style=python]
class Employee:
    def __init__(self, name):
        self.name = name

class Doctor(Employee):
    def __init__(self, name, specialty):
        super().__init__(name)
        self.specialty = specialty
\end{lstlisting}
  \begin{itemize}
    \item \texttt{super()} zpřístupní implementaci rodiče.
    \item Bez volání rodiče by atribut \texttt{name} nevznikl.
  \end{itemize}
\end{frame}

\begin{frame}{Dědičnost v UML}
  \centering
  \begin{tikzpicture}[>=Latex,node distance=17mm,
    cls/.style={draw,fill=SoftGray,minimum width=38mm,minimum height=12mm,align=center}]
    \node[cls] (parent) {\textbf{Employee}\\\texttt{name}};
    \node[cls,below=of parent] (child) {\textbf{Doctor}\\\texttt{specialty}};
    \draw[-{Triangle[open,length=4mm,width=5mm]}] (child)--(parent);
  \end{tikzpicture}
  \begin{itemize}
    \item Dutý trojúhelník míří od potomka k rodiči.
    \item Potomek dědí členy rodiče a přidává vlastní.
  \end{itemize}
\end{frame}

\begin{frame}[fragile]{Polymorfismus: stejné volání, různé chování}
\begin{lstlisting}[style=python,basicstyle=\ttfamily\normalsize]
class CashPayment:
    def process(self):
        return "Paid in cash"

class CardPayment:
    def process(self):
        return "Paid by card"

for payment in [CashPayment(), CardPayment()]:
    print(payment.process())
\end{lstlisting}
  Kód používá společné rozhraní a nemusí větvit podle konkrétního typu.
\end{frame}

% 0:38--0:42
\begin{frame}{Kompozice: objekt obsahuje jiné objekty}
  \begin{itemize}
    \item \texttt{Warehouse} uchovává seznam objektů \texttt{Item}.
    \item \texttt{Appointment} odkazuje na objekt \texttt{Patient} i \texttt{Doctor}.
    \item Kompozice vyjadřuje vztah „má“ (\emph{has-a}).
  \end{itemize}
  \begin{block}{Praktické pravidlo}
    Pro vztah „má“ preferujte kompozici; dědičnost použijte pro skutečnou specializaci „je druhem“.
  \end{block}
\end{frame}

\begin{frame}[fragile]{Speciální metody zapojují objekty do Pythonu}
\begin{lstlisting}[style=python]
class Vector2D:
    def __init__(self, x, y):
        self.x = x
        self.y = y

    def __str__(self):
        return f"({self.x}, {self.y})"

print(Vector2D(3, 4))    # (3, 4)
\end{lstlisting}
  Podobně \texttt{\_\_add\_\_} určuje chování operátoru \texttt{+} a \texttt{\_\_eq\_\_} porovnání \texttt{==}.
\end{frame}

\begin{frame}{Kdy třídu raději nevytvářet?}
  \begin{itemize}
    \item Jednorázový výpočet bez dlouhodobého stavu často stačí jako funkce.
    \item Třída pouze s jednou metodou bez atributů může být zbytečná vrstva.
    \item Nejprve se ptejte: existuje zde několik hodnot stejného druhu s vlastním stavem a chováním?
  \end{itemize}
  \begin{block}{Cíl není „použít OOP“}
    Cílem je kód, jehož struktura odpovídá problému a dobře se mění.
  \end{block}
\end{frame}

% 0:42--0:45
\begin{frame}{Časté chyby}
  \begin{itemize}
    \item Chybí dvojtečka za \texttt{class} nebo \texttt{def}.
    \item Metodě chybí první parametr \texttt{self}.
    \item Místo \texttt{self.name} se použije neexistující lokální \texttt{name}.
    \item Konstruktor se nejmenuje přesně \texttt{\_\_init\_\_}.
    \item Metoda je jen pojmenována: \texttt{account.deposit} místo volání se závorkami.
    \item Potomek přepíše \texttt{\_\_init\_\_}, ale neinicializuje rodičovské atributy.
  \end{itemize}
\end{frame}

\begin{frame}{Rychlá kontrola porozumění}
  \begin{enumerate}
    \item Jaký je rozdíl mezi třídou a objektem?
    \item Na co odkazuje \texttt{self} při \texttt{a.deposit(100)}?
    \item Proč mají dva účty rozdílný zůstatek?
    \item Co ukazuje diagram tříd a co sekvenční diagram?
    \item Co se stane, když potomek definuje metodu stejného jména jako rodič?
  \end{enumerate}
\end{frame}

\begin{frame}{Co využijete na cvičení?}
  \begin{itemize}
    \item \texttt{Rectangle}: atributy, metody a změna stavu.
    \item \texttt{Vector2D}: metody vracející nové objekty a \texttt{\_\_str\_\_}.
    \item \texttt{Warehouse} a \texttt{Item}: spolupráce objektů a kompozice.
    \item Hierarchie tvarů: dědičnost, přepis metod a polymorfismus.
    \item Pokročilé úlohy: další speciální metody a návrh více tříd.
  \end{itemize}
\end{frame}

\begin{frame}{Shrnutí}
  \begin{itemize}
    \item Třída je předpis, objekt je konkrétní instance s vlastním stavem.
    \item Atributy uchovávají data, metody definují chování.
    \item \texttt{self} označuje objekt aktuálního volání; \texttt{\_\_init\_\_} nastaví jeho počáteční stav.
    \item UML diagram tříd popisuje strukturu, sekvenční diagram průběh volání.
    \item Dědičnost specializuje, kompozice skládá a polymorfismus sjednocuje rozhraní.
  \end{itemize}
\end{frame}

\end{document}
